% ============================================================
% 1. Introduction
% ============================================================
\chapter{Introduction}

\section{Project Goal}

\textbf{CiboAmico} is a desktop application that supports a local food
marketplace: it lets a household monitor its pantry (\emph{frigo/dispensa}),
suggest compatible recipes based on what is actually available, and order fresh
products directly from local farmers and small suppliers. The system bridges
daily food management with short food-supply chains, reducing waste and
promoting local production.

The project is developed as the final assignment of the \emph{Ingegneria del
Software} (ISPW) course at Università di Roma ``Tor Vergata'', under the
supervision of Prof.\ Falessi and Prof.\ De Angelis (A.A.\ 2025/2026).

\section{Scope}

The system covers four functional areas:
\begin{itemize}
    \item \textbf{Pantry management}: add, list (ordered by expiry date) and
    monitor products in the household inventory;
    \item \textbf{Recipe engine}: find recipes compatible with the available
    ingredients, including an innovative ingredient-substitution mechanism;
    \item \textbf{Shopping list}: compute the missing ingredients for a chosen
    recipe;
    \item \textbf{Marketplace}: publish products as an approved local vendor,
    place single direct orders, track order state, and approve vendors and
    recipes as an administrator.
\end{itemize}

\section{Actors}

The system defines four primary actors:
\begin{itemize}
    \item \textbf{Client}: manages the pantry and uses the recipe
    engine and the marketplace to buy products;
    \item \textbf{Vendor}: a local farmer or supplier who, once
    approved by the administrator, publishes products with prices;
    \item \textbf{Nutritionist}: the only actor allowed to
    create recipes, which are published with a \emph{PROPOSED} state;
    \item \textbf{Administrator}: approves vendors and
    recipes.
\end{itemize}

Two external systems are used through adapters: \textbf{OpenFoodFacts}
(barcode-based product lookup) and the e-mail notification service
(\emph{Jakarta Mail}).

\section{Hardware and Software Requirements}

CiboAmico is a desktop application with the following minimum runtime
requirements:
\begin{itemize}
    \item Java Runtime Environment (JRE) 17 or later, with JavaFX 17+;
    \item 4~GB of RAM and 500~MB of free disk space;
    \item optional MySQL 8 server (the JDBC persistence mode); the
    file-system (JSON) and in-memory demo modes require no external
    services (NFR-01);
    \item a working e-mail account for the notification service
    (Jakarta Mail adapter) and an Internet connection for the
    OpenFoodFacts barcode lookup.
\end{itemize}

\section{Implementation Remarks}

The project deliberately implements the use cases under evaluation
(UC-01..UC-08, UC-11); UC-10 is documented for model completeness but is
excluded from the implementation and evaluation scope (see
Chapter~\ref{ch:usecases}).

\section{Document Structure}

The remainder of this document is organised as follows: Chapter~\ref{ch:related}
discusses related systems and the innovation elements of CiboAmico;
Chapter~\ref{ch:requirements} lists functional and non-functional requirements;
Chapter~\ref{ch:stories} presents the user stories (INVEST);
Chapter~\ref{ch:usecases} details the use cases with their internal steps;
Chapter~\ref{ch:storyboards} shows the interface storyboards;
Chapter~\ref{ch:design} describes the design (BCE architecture, VOPC, activity,
sequence and state diagrams); Chapter~\ref{ch:testing} presents the test plan
and results; Chapter~\ref{ch:exceptions} covers exception handling and
persistence; finally, Chapter~\ref{ch:links} collects the public links
(repository, SonarCloud report, demonstration video) and the conclusions.
